Class imbalance is often treated as the principal obstacle to rare-outcome prediction, yet adding minority observations cannot recover a decision boundary that is weak or poorly represented. This qualification examines whether class separability should be an explicit objective of synthetic tabular health-data generation. Early resampling, CTGAN augmentation, and causal-anchor filtering produced limited or inconsistent gains, motivating an explainability and neighbourhood analysis of class overlap. The resulting SepAware method generates an oversized CTGAN candidate pool and selects a class-count-preserving synthetic dataset according to realism, label separability, and domain-anchor compatibility. All reported predictive estimates use leakage-controlled five-fold stratified cross-validation on Brazilian Vigitel obesity data and a 334-record COVID-19 hospital mortality dataset. In a replicated $2^2$ experiment, separability accounts for 96.00\% of treatment-plus-error variation in Vigitel train-on-synthetic, test-on-real Macro-F1 ($p=5.03\times10^{-10}$), but only 5.13\% in COVID-19, where residual error accounts for 81.53\% and no treatment effect is significant. Anchor pressure is negligible in Vigitel and does not reliably improve COVID-19. These results support a dataset-dependent mechanism: explicit separability can recover useful class-conditional signal, but may also select easy prototypes and suppress clinically important overlap. The remaining PhD work therefore tests robustness across seeds, classifiers, generators, imbalance levels, and external datasets, while measuring minority diversity, causal consistency, clinical plausibility, and privacy. The intended contribution is not a universally superior generator, but a defensible framework for controlling and validating the separability--fidelity--utility trade-off in synthetic clinical data.
